% Terminal apparatus of the concise study: the scope note and the References
% for the sources this companion actually uses. The revision timestamp and the
% rights colophon follow from the shared generation record.

\section{Appendix: Scope and Qualifications}

{\footnotesize
\textbf{Edition, formulary and occurrence.} This concise study treats the Mass
\latin{Dominica decima nona post Pentecosten}, II classis, of the
\work{Missale Romanum}, editio typica, Typis Polyglottis Vaticanis 1962,
pp.~411--412, nos.~1679--1688, in the universal 1962 calendar and no particular
calendar. Its occurrence on 4~October 2026, the omission of the feast of St
Francis, the single oration of each kind, the colour, the Credo and the Trinity
Preface rest on the Missal's own calendar and general rubrics (RG 12, 16, 91,
111~b, 127~b; RGMR 433, 434~b, 475~a, 494~b). The optional external solemnity
of the Holy Rosary on the first Sunday of October rests on RGMR 358~b and 360;
the Rosary formulary it brings into up to two Masses is not studied here.

\textbf{Relation to the full study.} Every claim here is drawn from the full
study of this same formulary, the Full PDF subtitled \textit{The wedding
garment, the call to the highways, and the evening sacrifice: a study of the
proper in three readings}, which prints the appointed Latin with its public-domain English,
treats each element in its own setting, and develops the three readings at
length, each with its own four senses; nothing evidence-dependent is added here.
The four rows on the first page orient a reader and do not replace those
readings' own senses. The dossier's Date cells, relation labels, alternatives
and unresolved states are exactly those of the generated chronology record, as
in the study; the map and the dossier's explanatory rows condense the study's.
The joins between elements are the editor's; each Father and later writer is
cited only for what he said of his own passage, and Bl.\ Ildefonso Schuster for
the elements of this Mass he expounds.

\textbf{Text, English and rights.} Every Latin form of the ten elements was read
for the study on the page images of the CMAA facsimile of the typical edition
and compared with the Pustet \work{Missale Romanum} of Ratisbon, 1862, in which
each element's wording already stands; the collation is in
\texttt{propers/verified.md}. The Introit's antiphon was not traced to a
scriptural or liturgical source, and the Latin source of the Gradual's and
Offertory's added \latin{Domine} and of the Offertory's future tenses was not
identified. Psalms are cited in the Missal's Vulgate numbering; the map and the
dossier add the Hebrew numbering. The scriptural English is the Douay--Rheims in
Challoner's revision at the canonical verses, so that it does not render the
Missal's Latin where a chant departs from its verse; the English of the three
orations is the anonymous translation printed by Cummiskey in Philadelphia in
1861. Neither is an approved liturgical translation of the 1962 Missal. The
English of St Augustine's sermons and psalm expositions, of St John Chrysostom
on St Matthew and on Ephesians, and of St Irenaeus is that of the Nicene and
Post-Nicene and Ante-Nicene Fathers; English renderings of Latin not drawn from
a named translation are glosses of the Latin printed beside them. The Latin of
the formulary rests on its public-domain antecedent, the Pustet Missal (17
U.S.C. 103(b)); the Douay--Rheims, the 1861 English, Migne's texts, the
nineteenth-century translations of the Fathers, the Liège and Venice Aquinas,
O'Sullivan's 1866 Bellarmine and the 1927 English of Schuster's
\work{Sacramentary} are in the public domain in the United States. St Gregory's
Latin is quoted from the Latin Wikisource transcription of the
\work{Homiliarum in Evangelia}, licensed under Creative Commons
Attribution-ShareAlike 3.0, with attribution to its contributors at
\url{https://la.wikisource.org/wiki/Homiliarum_in_Evangelia}. The \work{New
American Bible} introductions are summarised, not reproduced.

\textbf{Reception and its limits.} The witnesses named here are those the
study's sweep read at the loci below. St Jerome, St Hilary and St Augustine's
\work{Quaestiones evangeliorum} were read on page images of Migne's
\work{Patrologia Latina}; St Gregory's homily in a transcription of Migne's text
not collated with the printed columns; St Thomas Aquinas on Ephesians on the
page images of the Liège edition of 1857, and on St Matthew only in the text
layer of the Venice edition of 1745, from which he is summarised and not
quoted. The Greek of St John Chrysostom on Pss~137 and 140 and of the
\work{Expositiones in Psalmos} printed under St Athanasius's name was read in
text layers of Migne's \work{Patrologia Graeca} and is described, not quoted.
Among the witnesses not reached are St Ambrose on Ps~118, Cassiodorus and
Theodoret on the psalms, the book of Origen's commentary on St Matthew that
treats this parable, the \work{Catena aurea}, St Anthony of Padua's sermon on
this Gospel and St Alphonsus Liguori's Sunday sermons.

\textbf{Chronology and records.} The Date cells on the second page come from
the Scripture chronology corpus through the record generated for this
formulary. That record gives the five psalms only the modern critical boundary
for the Psalter, from the introduction to the Psalms in the \work{New American
Bible, Revised Edition}, and holds no traditional date for them. For the Gospel
it holds five dates from the \work{Catholic Encyclopedia}'s article on St
Matthew and Durand's date for the Aramaic original; for the Epistle, Ladeuze's
range and the year of Prat's table. The articles do not agree, and the record
marks every one of these dates disputed. Beside them it carries two
comparisons requested under its critical profile, both from the \work{New
American Bible, Revised Edition}: a lower bound after A.D.~70 for the Greek
Gospel of Matthew, and around A.D.~80--100 for Ephesians, conditional on the
hypothesis of a later disciple; neither enters the default answer. The Gospel's
narrated event is undated.
The search, its loci, the disagreements and the negative results are in
\texttt{research/scope.md}, the reasoning behind the three readings in
\texttt{research/interpretations.md}, and this companion's production record in
\texttt{research/production-review.md}.
\par}

\section*{References}
\runninghead{References}

{\scriptsize
\subsection*{Liturgical books and Bibles}
\begin{itemize}\setlength{\itemsep}{0.15em}\setlength{\parskip}{0pt}
\item \work{Missale Romanum ex decreto Sacrosancti Concilii Tridentini restitutum, Summorum Pontificum cura recognitum}, editio typica (Typis Polyglottis Vaticanis, 1962), CMAA facsimile, \url{https://media.churchmusicassociation.org/pdf/missale62.pdf}: \work{Rubricae generales} 12, 16, 91, 111, 127; \work{Rubricae generales Missalis} 358, 360, 433, 434, 475, 494; \work{Ordo Missae} nos.~1032, 1037, 1137; pp.~411--412 (nos.~1679--1688).
\item \work{Missale Romanum} (Ratisbon: Pustet, 1862), pp.~351--352, \latin{Dominica XIX. post Pentecosten}.
\item \work{The Roman Missal, Translated into the English Language for the Use of the Laity} (Philadelphia: Eugene Cummiskey, 1861), pp.~445--448 (XIX Sunday after Pentecost, Collect, Secret, Postcommunion).
\item The Holy Bible, Douay--Rheims version, revised by Bishop Richard Challoner: 1 Par 15:3; 16:1, 7--8; Ps 77:1--8, 68--70; 104:1, 45; 118:4--5, 19, 54; 137:1--2, 7; 140:2; Mt 13:35; 21:1--18, 23, 43, 45; 22:1--15; 25:35; Rom 13:14; Eph 1:1; 3:1; 4:1, 22--32; 6:20; 1 Tim 1:5; 2:8. \work{Biblia Sacra Vulgatae editionis} (Clementine Vulgate): Ps 77:1; 104:1; 118:1, 4--5; 137:1, 7; 140:1--2.
\end{itemize}

\subsection*{Fathers, Doctors and saints}
\begin{itemize}\setlength{\itemsep}{0.15em}\setlength{\parskip}{0pt}
\item St Augustine, \work{Sermones} 90 and 95, Nicene and Post-Nicene Fathers, 1st ser., vol.~6: Sermo 90, 1, 4--6, 9; Sermo 95, 7. \work{Enarrationes in Psalmos}, 1st ser., vol.~8: on Ps~77, 2; 104, 1, 34; 118, Aleph 5; 137, 10; 140, 2--3; Latin on Ps~104, 1 and 34 in Migne's text (Latin Wikisource transcription). \work{Quaestiones evangeliorum} I.31, in Migne, PL 35 (Paris, 1845), col.~1329.
\item St Gregory the Great, \work{Homiliae in Evangelia} 38, 1, 3, 5--7, 9--14, Latin Wikisource transcription of Migne's text, CC BY-SA 3.0.
\item St Jerome, \work{Commentarii in Matthaeum} III, on Mt 22:1--14, and \work{Commentarii in Epistolam ad Ephesios} II, on Eph 4:23--28, in Migne, PL 26 (Paris, 1845), cols.~159--161, 508 and 512.
\item St Hilary of Poitiers, \work{Commentarius in Matthaeum} 22, 3--7, and \work{Tractatus super Psalmos}, in Ps.~118, Aleph 12, and in Ps.~140, 3--4, in Migne, PL 9 (Paris, 1844), cols.~1042--1044, 508--509 and 825--826.
\item St John Chrysostom, \work{Homilies on Matthew} 69 and \work{Homilies on Ephesians} 14, Nicene and Post-Nicene Fathers, 1st ser., vols.~10 and 13; \work{Expositio in Psalmos}, in Ps.~137 and 140, Migne, PG 55 (described, not quoted).
\item St Irenaeus of Lyons, \work{Against Heresies} IV.36.5--6, Ante-Nicene Fathers, vol.~1 (Buffalo, 1887), pp.~516--517.
\item \work{Expositiones in Psalmos} printed under the name of St Athanasius, in Ps.~118, Migne, PG 27 (described, not quoted).
\item St Thomas Aquinas, \work{In Epistolam ad Ephesios} c.~4, lect.~7--9 (Liège: H.~Dessain, 1857); \work{Super Evangelium S.~Matthaei lectura} c.~22 (Venice, 1745; summarised only).
\item St Robert Bellarmine, \work{A Commentary on the Book of Psalms}, trans. J. O'Sullivan (1866): on Pss~118:4; 140:2.
\item St Rabanus Maurus, \work{Commentariorum in Matthaeum libri VIII}, on Mt 22:8--10, in Migne, PL 107.
\end{itemize}

\subsection*{Liturgical commentator and chronology sources}
\begin{itemize}\setlength{\itemsep}{0.15em}\setlength{\parskip}{0pt}
\item Bl.\ Ildefonso Schuster, \work{The Sacramentary (Liber Sacramentorum)}, vol.~3 (London: Burns Oates \& Washbourne, 1927), pp.~172--174.
\item \work{The Catholic Encyclopedia} (New York, 1907--1912): ``Epistle to the Ephesians'' (vol.~5, 1909); ``Gospel of St.\ Matthew'' (vol.~10, 1911); ``St.\ Paul'' (vol.~11, 1911); ``The New Testament'' (vol.~14, 1912).
\item United States Conference of Catholic Bishops, \work{New American Bible, Revised Edition}, introductions to the Psalms, to Matthew and to Ephesians, on their dating; summarised, not reproduced.
\end{itemize}
\par}
